\scanpage{160}
\mspage{1 - 41 -}
\margin{n° 413}

Prenons maintenant $t\in U^{(s)}(G)$ : on a alors $t*\varepsilon_x\in T_x^{(s)}(X)$ et l'application $x\longrightarrow t*\varepsilon_x$ est une \uline{section} du fibré vectoriel $T^{(s)}(X)$ c'est-à-dire un \uline{opérateur différentiel} (au sens strict) \uline{d'ordre} $\leq s$ sur $X$. \uline{Nous le noterons souvent} $L_t^X$ \uline{dans la suite}, quand aucune confusion ne pourra en résulter. \uline{Cet opérateur est de classe} $C^{r-s}$ : en effet, l'application constante égale à $t$ est une application de classe $C^{r-s}$ de $X$ dans la variété $T^{(s)}(G)$ et l'application $x\longrightarrow\varepsilon_x$ est une section de classe $C^r$ du fibré $T^{(0)}(X)$ : notre assertion découle de là, puisque l'application $(u,v)\longrightarrow u*v$ est un morphisme bilinéaire de fibrés vectoriels de classe $C^{r-s}$ de $T^{(s)}(G)\times T^{(0)}(X)$ dans $T^{(s)}(X)$.

Si $t$ est sans terme constant, il en est de même de $t*\varepsilon_x$, donc de $L_t^X$.

En particulier, \uline{pour tout vecteur tangent} $t$, \uline{l'opérateur différentiel} $L_t^X$ \uline{est un champ de vecteurs} : c'est d'ailleurs le champ de vecteurs $x\longrightarrow tx$. \uline{L'application} $t\longrightarrow L_t^X$ \uline{de} $T_{\unsure{o|0}}(G)$ \uline{dans l'espace des champs de vecteurs sur} $X$ \uline{est appelée la loi d'opération infinitésimale associée à la loi d'opération donnée}.

PROPOSITION 23.- \uline{L'application} $t\longrightarrow L_t^X$ \uline{de} $U(G)$ \uline{dans l'espace des opérateurs différentiels sur} $X$ \uline{est linéaire}. \uline{Si} $t\in U^{(s)}(G)$, $t'\in U^{(s')}(G)$ \uline{avec} $s+s'\leq r$, \uline{on a} :
\[
L_{t'*t}^X=L_t^X\circ L_{t'}^X
\]
(\uline{au sens de la composition des opérateurs différentiels}).

La première assertion est évidente. Pour démontrer la seconde, considérons une fonction $\varphi$ de classe $C^r$ dans un ouvert $V$ de $X$ et soit $x\in V$. On a $L_{t'}\varphi(x)=t'(\varphi\circ\rho(x))$ et
\margin{159}
